\subsection{Galería de tiro}
Suponga que en la galería de tiro de un carnaval hay dos opciones para ganar
puntos:
\textbf{Opción 1:} Golpear un solo blanco, obteniendo el valor inscrito.
1
1
9
9
2
-5
-5
Puntaje +2
Opción 2: Golpear el disco que conecta dos consecutivos, obteniendo el pro-
ducto de los valores inscritos.
1
1
9
9
2
-5
-5
Puntaje +18
Tenga en cuenta: El puntaje inicial es cero y pueden haber blancos con
puntajes negativos. Está permitido no golpear un blanco para evitar perder
puntos.ISIS1105
PARCIAL 1 - 199320
3
2.1. [10 Pts]. Describa las entradas y salidas del problema.
E/S
Nombre
Tipo
Descripción
2.2. [20 Pts]. Defina una ecuación de recurrencia para L(k), el máximo puntaje
que se puede obtener jugando con los blancos 0, . . . , k - 1.
2.3. [10 Pts]. Dibujar el grafo de necesidades asociado con esta ecuación. Explique
qué representan los vértices y qué representan las aristas del grafo.
2.4. [20 Pts]. Desarrolle un algoritmo de programación dinámica para resolver
este problema.
